\documentclass[12pt]{ltjsarticle}
\usepackage{luatexja-fontspec}
\usepackage{amsmath,amssymb,mathtools}
\usepackage{physics}
\usepackage{tikz}
\usepackage{tikz-cd}
\usepackage{geometry}
\geometry{margin=25mm}

\title{線形代数課題 S3 詳解ノート}
\author{CODEX (OpenAI)でTeXファイルを生成\\CODEX (OpenAI)でLeanファイルを生成}
\date{2025年9月24日}

\begin{document}
\maketitle

\section{導入}
この文書は、ZEN大学「線形代数1」課題 S3 の各問題について、数学的背景と Lean 証明に対応する考え方を整理したものである。対象となる行列は全て実数体 $\mathbb{R}$ 上の行列とし、インデックス集合には有限集合 $\mathrm{Fin}(k)$ を用いる。Lean ファイル \texttt{S3.lean} に実装した証明がどのような定理に対応しているかを明瞭にし、順序対の集合と射（射影）による構造化の視点を意識した解説を与える。

\section{基礎問題の詳解}
以下、$A,B\in M_{m\times n}(\mathbb{R})$、$C\in M_{n\times p}(\mathbb{R})$、$x,y\in \mathbb{R}^n$、$z\in \mathbb{R}^m$ とする。

\subsection{S3\_P01 転置は加法に関して線形}
行列の転置は成分 $(i,j)$ に対し $A^{\mathsf{T}}_{ij}=A_{ji}$ を返す写像であり、加法写像と可換する。
\begin{align*}
  (A+B)^{\mathsf{T}}_{ij}
  = (A+B)_{ji}
  = A_{ji}+B_{ji}
  = (A^{\mathsf{T}}+B^{\mathsf{T}})_{ij}.
\end{align*}
したがって $(A+B)^{\mathsf{T}}=A^{\mathsf{T}}+B^{\mathsf{T}}$。Lean では `Matrix.transpose_add` がこの命題を直接表現しており、射影の観点では「転置」という反変関手が有限集合同士の直和に対してアーベル群準同型となる事実に相当する。

\subsection{S3\_P02 転置とスカラー倍の可換性}
実数 $c$ によるスカラー倍写像は各成分に $c$ を掛ける線形射である。したがって
\begin{align*}
  (c\cdot A)^{\mathsf{T}}_{ij} = (c\cdot A)_{ji} = c A_{ji} = (c\cdot A^{\mathsf{T}})_{ij}.
\end{align*}
Lean 補題 `Matrix.transpose_smul` を用いることで同様の等式を機械的に確認できる。

\subsection{S3\_P03 左分配法則}
行列積は中間インデックスでの有限和を用いて定義されるため、和に関して分配的である。
\begin{align*}
  \bigl(A(B+C)\bigr)_{ik}
  &= \sum_{j} A_{ij}(B_{jk}+C_{jk}) \\
  &= \sum_{j} A_{ij}B_{jk} + \sum_{j} A_{ij}C_{jk} \\
  &= (AB)_{ik} + (AC)_{ik}.
\end{align*}
Lean 側では `Matrix.mul_add` がこの性質を符号化している。これは $
abla$-diagram 上で、左からの射 $A$ が和射影を保存することを意味する。

\subsection{S3\_P04 積の転置}
転置は積を反転させる反変関手である。
\begin{align*}
  (AB)^{\mathsf{T}}_{ki} = (AB)_{ik} = \sum_j A_{ij}B_{jk}
  = \sum_j B^{\mathsf{T}}_{kj} A^{\mathsf{T}}_{ji} = (B^{\mathsf{T}}A^{\mathsf{T}})_{ki}.
\end{align*}
Lean では `Matrix.transpose_mul` によって実装している。これは圏論的には $\mathsf{T}$ が反変自動同型であることの言い換えである。

\subsection{S3\_P05 行列とベクトルの積の加法性}
列ベクトル $x,y$ に対する行列の左作用は線形で、有限和を通じて加法を保つ。
\begin{align*}
  A(x+y) = \left(\sum_j A_{ij}(x_j+y_j)\right)_i = (Ax)_i + (Ay)_i.
\end{align*}
Lean 補題 `Matrix.mulVec_add` は、この線形性をそのまま表現する。

\section{チャレンジ問題 S3\_CH の構造}
チャレンジ問題では内積 $\langle v,w\rangle = v^{\mathsf{T}}w$ を用い、転置が定義する随伴作用を示す。Lean 補題 `Matrix.dotProduct_mulVec` は $y^{\mathsf{T}}(Ax) = (A^{\mathsf{T}}y)^{\mathsf{T}}x$ を保証する。

\subsection{射影図式による理解}
以下の図式は、$A$ が $\mathbb{R}^n$ から $\mathbb{R}^m$ への線形射であり、転置 $A^{\mathsf{T}}$ が双対空間における随伴射であることを示す。
\begin{center}
  \begin{tikzcd}[column sep=large,row sep=large]
    \mathbb{R}^n \arrow[r, "A"] \arrow[d, "\mathrm{ev}_x"' description] &
    \mathbb{R}^m \arrow[d, "\mathrm{ev}_y" description] \\
    \mathbb{R} \arrow[r, equals] & \mathbb{R}
  \end{tikzcd}
\end{center}
ここで $\mathrm{ev}_x$ は列ベクトル $x$ による評価射、$\mathrm{ev}_y$ は $y$ による評価射である。図式可換性は
\begin{align*}
  \mathrm{ev}_y( A x ) = y^{\mathsf{T}}(Ax) = (A^{\mathsf{T}}y)^{\mathsf{T}} x = \mathrm{ev}_{A^{\mathsf{T}}y}(x)
\end{align*}
を意味し、Lean 証明では `Matrix.vecMul_transpose` による補助等式 `\mathsf{vecMul} \, y \, A = A^{\mathsf{T}} y` を用いて図式が可換になることを確認した。

\section{Lean 証明との対応表}
\begin{center}
  \begin{tabular}{lll}
    \hline
    問題 & 数学的内容 & 使用した Lean 補題/戦略 \\
    \hline
    S3\_P01 & 転置は加法写像 & \texttt{Matrix.transpose\_add} \\
    S3\_P02 & スカラー倍との交換 & \texttt{Matrix.transpose\_smul} \\
    S3\_P03 & 左分配法則 & \texttt{Matrix.mul\_add} \\
    S3\_P04 & 反変性 & \texttt{Matrix.transpose\_mul} \\
    S3\_P05 & 行作用の線形性 & \texttt{Matrix.mulVec\_add} \\
    S3\_CH & 随伴関係 & \texttt{Matrix.dotProduct\_mulVec}, \texttt{Matrix.vecMul\_transpose}, \texttt{dotProduct\_comm} \\
    \hline
  \end{tabular}
\end{center}

\section{まとめと今後の検討}
今回の課題では、行列の転置と内積が持つ反変性・随伴性が統一的に扱えることを確認した。特に、有限集合の積（順序対）上の射影・評価によって定義される内積は、転置が与える双対射によって整合的に構造化される。今後は、正定値内積空間や随伴行列（複素共役転置）への一般化、および Lean での自動証明支援（`simp?`, `aesop_math?`）の活用を検討するとよい。

\end{document}
